\section{Bốn nguyên nhân và cái giá của lối thoát}
\label{sec:ch15-nguyen-nhan}

\index{rò rỉ!bốn nguyên nhân|(}%
Mục trước để lại câu hỏi về chuyện gì xảy ra khi tập khái niệm thiếu. Có một
dụng cụ đo rồi thì câu ấy trả lời được bằng số, và bài báo trả lời cả nó lẫn ba
câu cùng loại: bốn quyết định thiết kế nào làm \tn{rò rỉ}{leakage} tăng, và tăng
bao nhiêu~\autocite{p30leakage}.

\index{rò rỉ!nguyên nhân thứ nhất, giám sát khái niệm quá nhẹ}%
Nguyên nhân thứ nhất là trọng số $\beta$ đặt quá thấp. Đây là nguyên nhân dễ
đoán nhất, vì $\beta$ thấp nghĩa là phép tối ưu ít bị phạt khi giá trị khái niệm
lệch khỏi nhãn khái niệm, nên nó tự do dùng cột ấy vào việc khác. Đo trên nhiều
bộ dữ liệu ở ba mức $\beta$, cả hai điểm số đều giảm khi $\beta$ tăng, đúng chiều
người ta chờ. Điều không ai chờ nằm ở hai đầu. Rò rỉ không về 0 khi $\beta$ tăng
mà chạm một mức sàn khác 0, và mức sàn ấy riêng cho từng lớp mô hình. Còn khi
$\beta$ lên quá cao, khoảng $\beta \ge 10$ trong các thí nghiệm của bài, mô hình
dồn hết vào việc học khái niệm và việc học nhiệm vụ hỏng, nên độ chính xác nhiệm
vụ tụt hoặc rò rỉ tăng trở lại. Mỗi lớp mô hình vì thế có một khoảng $\beta$ tối
ưu, và hai điểm số dùng để tìm khoảng ấy.

\index{rò rỉ!nguyên nhân thứ hai, mã hóa quá giàu}%
Nguyên nhân thứ hai là cách mã hóa giàu hơn mức cần. Khi nhãn khái niệm là bit
nhị phân mà giá trị học được là xác suất thực hay logit, phần dư thừa ấy là chỗ
trống, và phép tối ưu lấp nó. Bài báo phát biểu hai điều ở đây, và cả hai đều về
mô hình logit so với mô hình mềm. Mức sàn nói ở trên cao hơn hẳn với mô hình
logit. Và tăng $\beta$ kém hiệu quả hơn ở mô hình logit, nghĩa là không mua lại
được bằng giám sát cái mà cách mã hóa đã cho không. Cặp mô hình của
mục~\ref{sec:ch15-do-ro-ri} là một trường hợp cụ thể của cả hai: cùng bộ dữ
liệu, cùng mức $\beta$, chỉ khác cách mã hóa, và
\tn{điểm số can thiệp}{intervention score} lệch $0.301$.
Ở mức $\beta$ rất thấp thì khác biệt ấy biến mất, vì ở đó cách mã hóa nào cũng
rò rỉ nặng như nhau, không cách nào bị ràng buộc gì.

\index{rò rỉ!nguyên nhân thứ ba, tập khái niệm thiếu}%
Nguyên nhân thứ ba là tập khái niệm không đủ để giải nhiệm vụ. Khi bỏ bớt một
phần đáng kể các khái niệm được gán nhãn, thông tin cần cho nhiệm vụ không còn
nằm trọn trong tập ấy, nên mô hình cất phần thiếu vào những khái niệm còn lại.
Cách đo mức thiếu là huấn luyện riêng một đầu phân loại trên khái niệm thật rồi
xem nó đạt tới đâu: trên CUB, con số ấy là $1.000$ với tập đầy đủ 112 thuộc tính
và $0.836$ với tập đã rút gọn. Theo tiêu chí của
mục~\ref{sec:ch15-do-ro-ri}, rò rỉ tăng ở 14 trong 15 lớp mô hình huấn luyện
trên tập khái niệm thiếu, và chênh lệch rõ nhất ở các mức $\beta$ thấp và trung
bình.

\index{rò rỉ!nguyên nhân thứ tư, đầu phân loại quá nghèo}%
Nguyên nhân thứ tư đối xứng với nguyên nhân thứ ba và ít hiển nhiên hơn: đầu
phân loại quá nghèo để học quan hệ giữa khái niệm và nhãn. Chỗ khó xử là đầu
phân loại tuyến tính lại chính là lựa chọn mà người ta chuộng, vì một tầng tuyến
tính trên bảng khái niệm cho một lời giải thích đọc thẳng ra được từ trọng số.
Nhưng nếu nhãn thật là một hàm phi tuyến của các khái niệm thì tầng ấy không
biểu diễn nổi, và phép tối ưu bù bằng cách bóp méo chính các giá trị khái niệm
cho tới khi một hàm tuyến tính của chúng khớp được nhãn. Đo trên các bộ dữ liệu
đã sửa nhãn thành hàm phi tuyến của khái niệm, rò rỉ tăng ở 7 trong 9 lớp mô
hình, và tăng $\beta$ nói chung không bù lại được.

\index{rò rỉ!bốn nguyên nhân|)}%
Bốn nguyên nhân ấy không nằm cùng một hạng với nhau, và
mục~\ref{sec:ch15-phan-con-lai} sẽ thấy khảo sát của ngành xếp chúng khác. Hai
nguyên nhân đầu là \tn{siêu tham số}{hyperparameter}: người dựng mô hình chọn
được, và hai điểm số
chỉ ra chọn thế nào. Hai nguyên nhân sau là chuyện dữ liệu và chuyện nhiệm vụ,
và không siêu tham số nào gỡ được. Đó là chỗ mà cái giá của lối thoát bắt đầu
hiện ra.

\index{CEM!rò rỉ theo thiết kế|(}%
Cái giá ấy rõ nhất ở kiến trúc được coi là bước tiến của họ mô hình này. CEM,
tức mô hình nhúng khái niệm, thay mỗi giá trị khái niệm bằng một cặp vector, một
cho trạng thái có mặt và một cho trạng thái vắng mặt, rồi trộn hai vector theo
xác suất khái niệm và đưa vector trộn vào đầu phân loại. Lý do người ta dựng nó
là hiệu năng: CEM đạt độ chính xác nhiệm vụ cao hơn hẳn mô hình nút thắt thông
thường, thường ngang mô hình đầu cuối không có tầng thắt nào, và vì thế được xem
là đã gỡ được thế đánh đổi giữa độ chính xác và tính diễn giải được.

\index{CEM!thông tin đi tự do từ đầu vào tới nhãn}%
Bài báo đọc nó theo chiều ngược, và lập luận nằm ở chỗ giám sát đặt vào đâu.
Trong CEM, mất mát trên khái niệm chỉ tác động lên xác suất khái niệm, còn thứ
thật sự đi vào đầu phân loại là vector trộn, và không có mất mát nào đặt lên
vector ấy. Nên trong lượt truyền xuôi, một lượng thông tin tùy ý được phép chảy
từ đầu vào tới dự đoán nhãn, miễn là hai vector sinh ra vẫn đủ dự báo được cái
bit khái niệm. Nói cách khác,
\tn{rò rỉ khái niệm-nhiệm vụ}{concepts-task leakage} ở đây không phải một tai
nạn của phép tối ưu mà là thứ kiến trúc chừa sẵn chỗ cho. Đo thì thấy đúng vậy:
lượng thông tin mà vector trộn mang về nhãn đi cùng chiều với độ chính xác nhiệm
vụ, trong khi độ chính xác khái niệm thì không, nghĩa là hiệu năng của CEM đến
từ phần thông tin thêm chứ không từ phần khái niệm. Cùng hướng với điều ấy là một
kết quả mà chính các tác giả CEM báo cáo: bỏ tới 90 phần trăm số khái niệm được
gán nhãn của CUB, độ chính xác nhiệm vụ của CEM gần như không đổi.

\index{CEM!rò rỉ liên khái niệm tăng theo giám sát}%
\tn{Rò rỉ liên khái niệm}{interconcept leakage} ở CEM còn đi ngược quy luật của mục trên. Ở mô hình nút thắt
thông thường, tăng $\beta$ làm cả hai điểm số giảm. Ở CEM, tăng giám sát khái
niệm làm vector trộn dự báo khái niệm của chính nó tốt hơn, đúng như mong muốn,
nhưng đồng thời làm nó dự báo các khái niệm khác cũng tốt hơn, tức rò rỉ liên
khái niệm tăng. Đo trên xác suất khái niệm thì hai điểm số vẫn giảm như ở mô hình
thường; đo trên vector trộn thì không. Vậy thông tin rò rỉ
không mất đi mà chuyển chỗ, từ xác suất khái niệm sang vector, và vector là chỗ
mà đầu phân loại thật sự đọc. Nên với CEM không có giá trị siêu tham số nào vừa
cải thiện việc học khái niệm vừa giảm rò rỉ liên khái niệm.

\index{CEM!rò rỉ căn chỉnh}%
Bài báo còn tìm ra một dạng rò rỉ thứ ba, đặt tên là
\tn{rò rỉ căn chỉnh}{alignment leakage}, và nó chỉ
xuất hiện ở mô hình được huấn luyện kèm can thiệp. CEM trong lúc huấn luyện thỉnh
thoảng bị đặt xác suất khái niệm về giá trị thật, để nó quen với việc bị can
thiệp. Mô hình có thể học cách lợi dụng điều đó: sinh ra vector \enquote{đúng
phía} dự báo nhãn tốt hơn vector phía kia, nên can thiệp làm hiệu năng tăng mà
không phải vì khái niệm được sửa đúng. Đo trên năm bộ dữ liệu thì hiện tượng ấy
vắng mặt ở ba và có mặt rõ ở hai, và ở hai bộ có nó thì nó mạnh hơn khi giám sát
khái niệm và tần suất can thiệp lúc huấn luyện đều cao. Hệ quả cho mục~\ref{sec:ch15-dung-cu-duoc-kiem}
là hệ quả đã nói ở đó: với kiến trúc kiểu này thì hiệu năng khi can thiệp thôi
không còn đo rò rỉ nữa.

\index{CEM!rò rỉ theo thiết kế|)}%
\index{mô hình nút thắt khái niệm!cái giá của lối thoát}%
Mục 9 của bài báo gom bốn nguyên nhân thành một quy trình thiết kế, và quy
trình ấy khác các quy trình cùng loại ở chỗ nó chịu dừng. Bước đầu là kiểm xem tập khái
niệm có đủ dự báo nhiệm vụ không, bằng cách huấn luyện vài đầu phân loại trên
khái niệm thật, bắt đầu từ tầng tuyến tính. Nếu ngay cả kiến trúc phức tạp tùy ý
cũng không đạt độ chính xác đủ cao thì tập khái niệm thiếu, và bài báo viết rằng
khi không kiếm thêm được nhãn khái niệm thì phải kết luận nhiệm vụ ấy không phù
hợp với lối tiếp cận theo khái niệm trên bộ dữ liệu đang có. Câu ấy là cái giá
của lối thoát viết ra thành một dòng, và nó do chính người đề xuất khung viết.
Các bước sau chọn cách mã hóa không giàu hơn mức cần, rồi chọn $\beta$ trong
khoảng tối ưu; và bài báo ghi rằng ngay cả khi làm hết như vậy, rò rỉ ở mô hình
mềm thường vẫn không tránh được, chỉ là ở mức mà mô hình dựa vào nó ít hơn dựa
vào giá trị khái niệm.
